\section{Zero loss over a fixed coefficient field}\label{sec:field}
\noindent
Suppose all the polynomials are defined over a finite extension
$K/\mathbf Q_p$ of absolute ramification index $e$. For an irreducible
factor of degree $d$ containing a root of valuation $1/N_h$, all its roots
have that valuation, by invariance of the unique extended valuation.
The norm has valuation $d/N_h\in e^{-1}\mathbf Z$. A nonempty shell thus
costs at least $1/\gcd(e,N_h)\ge1/e$. Among any $\lfloor eB\rfloor+1$
eligible shells there is one with no root at its radius. All primitive
evaluations there are simultaneously bounded above by their witness
valuations. 


\subsection*{An exact formula over a cyclotomic coefficient field}
Fix $s\ge1$, a compatible system of primitive roots $\xi_h$ with $\xi_{h+1}^p=\xi_h$, and the coefficient field $K_s=\mathbf Q_p(\xi_s)$. Its absolute ramification index is $N_s$. On an arbitrarily high shell, a packet over this field has precisely the same optimum as a packet on shell $s$.

\begin{theorem}[Descent to the coefficient level]\label{thm:descent}
For packets with coefficients in $\mathcal O_{K_s}$, the exact minimum witness budget needed to block the entire interval $[h_0,H]$ at tolerance $\delta\ge0$ is
\[
       \sum_{h=h_0}^H a_{\min(h,s)}(\delta),
       \qquad a_j(\delta)=m_j/N_j .
\]
For $h_0\ge s$ the least guaranteed order exponent is exactly
\[
        h_0+\left\lfloor B/a_s(\delta)\right\rfloor .
\]
\end{theorem}
\begin{proof}
For $h\le s$ both the lower bound and the optimizing linear factors of Theorem~\ref{thm:exact-obstruction} already apply. Suppose $h>s$. Partition all roots of valuation $r_h$ in the packet into conjugacy orbits over $K_s$, repeated according to multiplicity. If an orbit has degree $d$, its norm valuation $d/N_h$ lies in $N_s^{-1}\mathbf Z$, so $u=dN_s/N_h$ is a positive integer.

Pick a member $\alpha$ and a nearest point $\xi_h^b-1$. Every automorphism fixing $K_s$ acts on $\xi_h$ by an exponent congruent to $1$ modulo $p^s$. Consequently each root in that orbit has a nearest point whose exponent is congruent to $b$ modulo $p^s$. Ties cause no problem: one conjugating automorphism can be chosen for each member. Assign $u$ marks to the single unit class $b\bmod p^s$. If a test class $a\bmod p^s$ is unmarked, its $p$-adic distance in exponent from each of the nearest exponents for this orbit is given by the same $k=v_p(a-b)<s$. At every lift of $a$ to shell $h$, this orbit contributes to the upper bound for the loss at most
\[
             d\,\frac{p^k-1}{N_h}=u\,\frac{p^k-1}{N_s}.
\]
Write $U=\sum u$. The actual root budget on shell $h$ is $\sum d/N_h=U/N_s$. If $U<m_s\le N_s$, the least-occupied path on shell $s$ supplies an unmarked class with total loss at most $S_s(U)\le\delta$ at every lift. Thus blocking requires a cost at least $a_s$ on each shell $h>s$.

Conversely choose the first $m_s$ positive unit exponents $b$ at level $s$. At level $h>s$ take one polynomial for each such $b$,
\[
            F_{h,b}(T)=(1+T)^{p^{h-s}}-\xi_s^b.
\]
Its roots are exactly the $\xi_h^c-1$ with $c\equiv b\pmod {p^s}$. Its value at the origin and its Gauss baseline both have valuation $r_s$. At $T=\xi_h^a-1$, its value is $\xi_s^a-\xi_s^b$. Hence the entire loss profile at level $s$ pulls back to level $h$. Values on all other shells have zero positive excess by the root factorization. Combining these polynomials with the linear constructions below level $s$ realizes the formula with a common witness $0$. The last assertion follows from strict budget crossing.
\end{proof}
For example, take $p=2$, $s=3$ and $\delta=1/4$. At every $h\ge3$ the exact price of blocking the shell over $\mathbf Q_2(\zeta_8)$ is $a_3=1/2$. A field that already contains all $2^h$-th roots permits the smaller costs in the preceding table as $h$ grows. This identifies exactly where extra coefficient ramification can aid an adversarial packet.


\subsection*{Several variables}
The fixed-field result does not require witnesses in the coefficient field. Let $K/\mathbf Q_p$ be finite of absolute ramification index $e$ and take a finite family $F_i\in\mathcal O_K[T_1,\dots,T_s]\setminus\{0\}$. Write
\[
 A_i=\min_{\mathbf a}\{v([T^{\mathbf a}]F_i)+\rho|\mathbf a|\},
 \qquad \sum_i A_i\le B .
\]
Then there is a single tuple $z_j=\zeta_j-1$, each of order exponent $h_j\ge h_0$ for $\zeta_j$, such that
\[
                 v(F_i(z_1,\ldots,z_s))\le A_i\quad\text{for every }i.
\]
This gives zero positive excess relative to all witnesses of coordinate valuations at least $\rho$. Explicit exponents $h_j\le M_j$ are available before the polynomials are given:
\[
 M_1=h_0+\lfloor eB\rfloor,\qquad
 M_{j+1}=h_0+\lfloor e N_{M_j}B\rfloor.
\]

\begin{proof}
At a given step expand a partially specialized polynomial as
$\sum_{\nu} g_{i,\nu}(T)Y^\nu$, where $Y$ denotes all later variables. Its multivariate Gauss valuation $A_i'$ is $\min_\nu(G_\rho(g_{i,\nu})+\rho|\nu|)$. Choose one nonzero coefficient polynomial $g_i$ attaining that minimum. Then $\sum_i G_\rho(g_i)\le\sum_i A_i'\le B$, initially by hypothesis. Apply the one-variable fixed-field result with these Gauss baselines, which the root-budget proof also permits. At a common point $z$ one has $v(g_i(z))\le G_\rho(g_i)$ for every $i$. The new Gauss valuation after specialization is at most $v(g_i(z))+\rho|\nu|\le A_i'$. This also ensures every new polynomial is nonzero and integral. Iterate over all coordinates.

The coefficient field after $j$ steps is contained in $K(\mu_{p^{M_j}})$, as the stated bounds are nondecreasing and the cyclotomic levels are nested. Its absolute ramification index is at most $e[K(\mu_{p^{M_j}}):K]\le eN_{M_j}$. The one-variable count of root-bearing shells gives the recurrence. At the final coordinate the Gauss valuation of the resulting constant is its valuation.
\end{proof}

\paragraph{What this gives for arithmetic twist families.}
This bound can be used whenever the coefficient field is controlled \emph{before} choosing the order cutoff. The worst-case packets in Theorem~\ref{thm:exact-obstruction} live in finite cyclotomic fields whose ramification grows with the largest forbidden shell. An application to a particular Selmer or special-value family therefore turns on its coefficient fields, or on structural restrictions on its interpolation polynomials. The extremal packets are explicit integral polynomials; their occurrence as Selmer determinants would be an additional arithmetic theorem.

\paragraph{Further research questions.}
Determine the optimum at fixed packet cardinality between the single-polynomial and unrestricted regimes. Determine the exact obstruction cost over a finite extension that is not cyclotomic. In several variables, is the iterated growth of the displayed fixed-field bound necessary, or can one use the common Galois action on the coordinates to reduce the search? The independent analytic applications in \texttt{prime-sampling.pdf} pursue zero location and simplicity by a different part of the same Lean proof.
